\chapter*{Abstract}
\addcontentsline{toc}{chapter}{Abstract}

Meeting a timing constraint is as much a requirement of a correct digital
circuit as its logical function, yet timing correctness cannot be verified by
functional simulation alone. This report documents the design,
implementation, and evaluation of \tool, a software tool that performs
static timing analysis (STA) on a flattened gate-level netlist and then
improves an initial, possibly timing-violating, gate sizing using a
logical-effort--guided optimizer. The tool ingests a netlist, a sized
standard-cell library, and a configuration file; validates the netlist's
structural correctness; builds a combinational directed acyclic graph (DAG);
performs rise/fall--separated forward and backward timing propagation under
an RC-based linear delay model with unate-aware timing arcs; extracts the
top-$K$ least-slack critical paths; performs a logical-effort decomposition
of each path ($g$, $h$, $b$, $G$, $H$, $B$, $F$, and the theoretical minimum
path delay); estimates leakage and activity-driven dynamic power together
with total area; and finally searches for a gate-sizing assignment that
removes timing violations while respecting area and power budgets, using a
cost function that trades off delay, power, area, and worst-negative-slack
violation.

Four gate-sizing heuristics are implemented and compared head-to-head under
identical stopping and acceptance rules: an exhaustive logical-effort-guided
search, two logical-effort-informed ranking heuristics of differing
aggressiveness, and a logical-effort-free random-greedy baseline. Beyond the
scope required by the course assignment, the project additionally includes
(i) a synthetic benchmark generation and evaluation framework that produces
statistically meaningful sizing-repair problems with a certified best-known
reference solution and evaluates every heuristic's decisions against an
independent counterfactual oracle, and (ii) an interactive, browser-based
topology viewer for visually inspecting a circuit's pre- and
post-optimization state. A Monte Carlo statistical timing analysis models
global and local manufacturing process variation and reports the resulting
delay distribution, timing yield, and critical-path stability before and
after optimization.

The tool is validated against the fifteen valid and six intentionally
invalid benchmark circuits supplied with the assignment. Results show that
logical-effort-guided sizing resolves a substantially larger fraction of
timing violations than a logical-effort-free random-greedy baseline
(13 of 14 originally-violating circuits versus 7 of 14), reaches a final
design cost that is never worse and strictly lower on two-thirds of the
suite, and does so with fewer total static-timing-analysis calls summed
across the whole suite. The synthetic benchmark suite provides a statistically
grounded confirmation of this result across a much larger and more diverse
population of circuits than the fixed benchmark set alone could support.

\vfill
\noindent\textbf{Keywords:} static timing analysis, logical effort, gate
sizing, standard-cell library, VLSI CAD, timing closure, Monte Carlo
statistical timing analysis, benchmark generation.

\clearpage